% p1-11-reviewer-risk

\section{P1 §11. Reviewer-Risk Assessment}

11. Reviewer-Risk Assessment

      Figure (fig-p1-ch11-reviewer). Reviewer-path triptych --- R1..R6 risk cards / response arrows /
                                         honest-admission seal.

  11.1 Purpose
  This section anticipates the principal objections that a journal reviewer with expertise
  in information theory, particle physics, or Bayesian statistics would raise. It documents
  the risks honestly and explains how the methodological choices address them.

  11.2 Risk R1: Multiple Testing Inflation
  Risk: The hypothesis class S(C) is large enough that small residuals are guaranteed
  even for random targets. The Bonferroni correction may be insufficient if grammar
  values are correlated.

  Mitigation: The effective trial count Neff(C) = |S(C)| is used as the Bonferroni
  denominator, which is conservative when grammar values are positively correlated
  (closely spaced lattice points). The permutation calibration of Section 6.5 provides an
  empirical calibration against the actual null distribution, accounting for correlations
  automatically.

  11.3 Risk R2: Grammar Density Masquerading as Selectivity
  Risk: The $\varphi$-lattice may be denser in the empirical range of physical constants than in
  the null model's range, creating apparent compressibility that is purely geometric.

  Mitigation: The Mperm null model preserves the empirical range of T and directly
  controls for this. The Mrand null model additionally tests whether the specific algebraic
  structure of the basis, as opposed to its cardinality, is responsible for observed
  coverage. The negative controls of Section 10.5 provide a direct diagnostic.

  11.4 Risk R3: Post-Hoc Data Selection

  Risk: The target dataset T was assembled after observing fit results, introducing
  selection bias.

  Mitigation: Protocol P1 (Section 4.3) mandates pre-registration of T before any
  symbolic search is executed, with a cryptographic hash commitment. The dataset is
  defined by membership criteria (dimensionless, Standard Model, PDG-tabulated,
  precision > 10-2), not by fit quality.

  11.5 Risk R4: Non-Uniqueness of the Grammar Basis
  Risk: The choice of basis {$\varphi$, pi, e, 3} is arbitrary; a different basis might yield similar
  or better results by accident.

  Mitigation: The paper explicitly claims no uniqueness for the basis. The Mrand null
  model tests whether results are specific to the structured basis. Future work should
  compare G$\varphi$ against alternative algebraically structured grammars using the same
  MDL framework.

  11.6 Risk R5: Misinterpretation as Derivation
  Risk: Readers may interpret low-residual fits as derivations of physical constants or as
  evidence of an underlying physics principle.

  Mitigation: Section 1.3 explicitly disavows all physical derivation claims. The encoding
  map is repeatedly characterized as a constrained quantization, not a physical law.
  Every expression is accompanied by the remark that it is a property of the
  representation system, not a property of nature.

  11.7 Risk R6: PhD Monograph Provenance Bias
  Risk: The paper's methodological design and candidate expressions may be biased by
  the PhD monograph's exploratory phase.

  Mitigation: Section 12 provides a transparent delineation of what the monograph
  contributes (algebraic structure, falsification design, historical motivation) versus what
  it does not provide (statistical evidence). All expressions from the monograph's Chapter
  34 catalogue are subject to the same null-model analysis as expressions not in the
  catalogue.

